\section{Model}

We test three nested objects. First, the original CGHSTC-33-RF core defines a
calibration clock factor \(\Gamma_{\rm cal}\) and a propagation clock factor
\(\Gamma_{\rm prop}(z)\) from a shared coupling and selector structure. Second,
H0-CDR defines a domain-table approximation that maps observed calibration
methods onto predeclared clock-domain selectors. Third, method-split candidates
such as Cepheid-vs-population indicators are diagnostic baselines, not the full
CGHSTC-33-RF theory.

The calibration-domain response model applies:

\begin{equation}
H_{0,i}^{\text{pred}} = \frac{H_0^{\text{cos}}}{\Gamma_{\text{cal}}(d_i)}
\end{equation}

where $\Gamma_{\text{cal}}(d_i) = 1 - \epsilon_{\text{cal}} \, S_{\text{cal}}(d_i)$
and $S_{\text{cal}}$ is computed from pre-declared domain features.

The route-dependent coupling form is retained as an exploratory extension only.
In the present sample-placeholder setting it is not treated as evidence for a
physical route-dependent epsilon term.

The propagation-clock factor modifies the comoving distance:

\begin{equation}
D_C(z) = c \int_0^z \frac{\Gamma_{\text{prop}}(z')}{H_{\text{LCDM}}(z')} \, dz'
\end{equation}
